% ===================== INTRODUCCION =====================
\section{Introduccion}
\label{sec:intro}

Las tuber\'{i}as de proceso de una planta industrial operan sometidas
simult\'{a}neamente a presi\'{o}n interna, peso propio y gradientes t\'{e}rmicos
entre la condici\'{o}n de instalaci\'{o}n y la de operaci\'{o}n. La expansi\'{o}n
t\'{e}rmica restringida por anclas, equipos y soportes induce esfuerzos y cargas
que, de no verificarse, pueden comprometer la integridad de la tuber\'{i}a, de las
boquillas de los equipos conectados y de los propios soportes. El an\'{a}lisis de
flexibilidad es la verificaci\'{o}n reglamentaria que el c\'{o}digo ASME B31.3
exige para demostrar que un sistema de tuber\'{i}as absorbe dichas deformaciones
dentro de l\'{i}mites seguros~\cite{asme_b313_2016}. El presente informe documenta
el an\'{a}lisis de flexibilidad de la l\'{i}nea de succi\'{o}n del tanque TK
Blowtank (SIM-012) del proyecto P2603 SW-K60, ejecutado para Smurfit Westrock
con el software CAESAR II 2019~\cite{caesar2_ug_2019}.

% ─────────────────────────────────────────────────────────────────────────────
\subsection{Antecedentes}

El proyecto P2603 SW-K60 comprende la modificaci\'{o}n y ampliaci\'{o}n de
sistemas de tuber\'{i}as de la planta de Cart\'{o}n Colombia de Smurfit Westrock.
Dentro del paquete de an\'{a}lisis de flexibilidad se han evaluado ocho
l\'{i}neas asociadas a los sistemas de lavado (DDW), recirculaci\'{o}n, vac\'{i}o
y transferencia; todas ellas comparten las mismas bases de dise\~{n}o
(secci\'{o}n~\ref{sec:bases}). La l\'{i}nea objeto de este informe corresponde a
la succi\'{o}n desde el tanque TK Blowtank, modelada a partir del isom\'{e}trico
PCF111 emitido por el equipo de dise\~{n}o en AutoCAD Plant 3D 2027.

% ─────────────────────────────────────────────────────────────────────────────
\subsection{Descripcion del Sistema}

La l\'{i}nea SIM-012 conduce agua desde la boquilla de succi\'{o}n del tanque TK
Blowtank hacia el sistema de bombeo aguas abajo. Est\'{a} construida en acero
inoxidable ASTM A312 TP 304L, predominantemente Sch 10S con tramos Sch 40S, en
di\'{a}metros nominales de 6 a 16~NPS, e incluye codos de 90$^\circ$ y
45$^\circ$ de radio largo (ASME B16.9), reducciones exc\'{e}ntricas, uniones
bridadas tipo \textit{lap joint} con \textit{stub-end} (ASME B16.5) y v\'{a}lvulas
de compuerta de cuchilla. El modelo de elementos finitos comprende 58 nodos
evaluados, con restricciones tipo soporte r\'{i}gido vertical (+Y) y dos anclas
(ANC) en los extremos del tramo analizado.

% ─────────────────────────────────────────────────────────────────────────────
\subsection{Motivacion del Estudio}

Toda l\'{i}nea nueva o modificada del proyecto debe demostrar el cumplimiento de
los criterios de esfuerzo de ASME B31.3-2016 antes de su puesta en servicio.
La pregunta t\'{e}cnica que responde este informe es si la configuraci\'{o}n
geom\'{e}trica y de soporter\'{i}a de la l\'{i}nea SIM-012 mantiene los esfuerzos
de c\'{o}digo, los desplazamientos y las cargas en soportes dentro de l\'{i}mites
aceptables en todas las condiciones de operaci\'{o}n previstas.

% ─────────────────────────────────────────────────────────────────────────────
\subsection{Marco Teorico}
\label{ssec:marco}

\subsubsection{Expansion termica y esfuerzo en tuberias restringidas}

Una tuber\'{i}a instalada a temperatura ambiente y operada a mayor temperatura
experimenta una expansi\'{o}n axial libre dada por:
\begin{equation}
    \Delta L = \alpha \, L \, \Delta T,
    \label{eq:expansion}
\end{equation}
donde $\alpha$ es el coeficiente de expansi\'{o}n t\'{e}rmica lineal del material,
$L$ la longitud del tramo y $\Delta T$ la diferencia entre la temperatura de
operaci\'{o}n y la de instalaci\'{o}n. Si los extremos del tramo estuvieran
totalmente restringidos, la expansi\'{o}n se convertir\'{i}a en un esfuerzo de
compresi\'{o}n de magnitud:
\begin{equation}
    \sigma = E \, \alpha \, \Delta T,
    \label{eq:esfuerzo_termico}
\end{equation}
con $E$ el m\'{o}dulo de elasticidad del material~\cite{toolbox_thermal_stress}. Para aceros inoxidables
austen\'{i}ticos y los $\Delta T$ t\'{i}picos de proceso, este esfuerzo superar\'{i}a
con facilidad el admisible; por ello las configuraciones reales incorporan codos
y cambios de direcci\'{o}n que aportan flexibilidad y reducen las cargas
transmitidas a anclas y equipos~\cite{peng_pipe_stress}.

\subsubsection{Modelo numerico: metodo de elementos finitos}

CAESAR II resuelve el sistema de tuber\'{i}as mediante el m\'{e}todo de elementos
finitos: la l\'{i}nea se discretiza en elementos viga tridimensionales con seis
grados de libertad por nodo (tres traslaciones y tres rotaciones), y el equilibrio
est\'{a}tico del sistema se expresa en forma matricial:
\begin{equation}
    \mathbf{K} \, \mathbf{u} = \mathbf{F},
    \label{eq:fem}
\end{equation}
donde $\mathbf{K}$ es la matriz de rigidez global ensamblada a partir de la
geometr\'{i}a, el material y las restricciones de soporte, $\mathbf{u}$ el vector
de desplazamientos nodales y $\mathbf{F}$ el vector de cargas (peso, presi\'{o}n,
desplazamientos t\'{e}rmicos impuestos). Resueltos los desplazamientos, el
programa calcula los esfuerzos en cada elemento y los compara con los
admisibles del c\'{o}digo seleccionado~\cite{caesar2_ug_2019}.

\subsubsection{Criterios de esfuerzo de ASME B31.3-2016}

El c\'{o}digo distingue dos familias de cargas. Las \textbf{cargas sostenidas}
(peso y presi\'{o}n) son primarias, de car\'{a}cter no auto-limitante: su esfuerzo
longitudinal $S_L$ (evaluado seg\'{u}n el p\'{a}rrafo 320.2) no puede exceder el
admisible b\'{a}sico a temperatura de operaci\'{o}n (p\'{a}rrafo 302.3.5(c)):
\begin{equation}
    S_L \leq S_h.
    \label{eq:sl}
\end{equation}

Las \textbf{cargas t\'{e}rmicas} son secundarias y auto-limitantes: una fluencia
local redistribuye y relaja el esfuerzo, y el modo de falla asociado es la fatiga
por ciclos. El c\'{o}digo eval\'{u}a el rango de esfuerzo de expansi\'{o}n $S_E$
(p\'{a}rrafo 319.4.4, ecuaci\'{o}n (17)):
\begin{equation}
    S_E = \sqrt{S_b^2 + 4 S_t^2},
    \label{eq:se}
\end{equation}
donde $S_b$ es el esfuerzo de flexi\'{o}n resultante, amplificado por los
factores de intensificaci\'{o}n de esfuerzo $i_i$ (en el plano) e $i_o$ (fuera del
plano) propios de cada accesorio (Ap\'{e}ndice D de B31.3-2016):
\begin{equation}
    S_b = \frac{\sqrt{(i_i M_i)^2 + (i_o M_o)^2}}{Z},
    \label{eq:sb}
\end{equation}
y $S_t$ el esfuerzo torsional:
\begin{equation}
    S_t = \frac{M_t}{2 Z},
    \label{eq:st}
\end{equation}
con $M_i$, $M_o$ y $M_t$ los momentos flectores en el plano, fuera del plano y
torsor, y $Z$ el m\'{o}dulo resistente de la secci\'{o}n. El criterio de
aceptaci\'{o}n es (p\'{a}rrafo 302.3.5(d), ecuaci\'{o}n (1a)):
\begin{equation}
    S_E \leq S_A = f \left( 1.25 \, S_c + 0.25 \, S_h \right),
    \label{eq:sa}
\end{equation}
con $f$ el factor de reducci\'{o}n por n\'{u}mero de ciclos t\'{e}rmicos previstos
(Tabla 302.3.5; $f = 1.0$ hasta 7000 ciclos). Cuando $S_h > S_L$, el c\'{o}digo
permite la forma liberal de la ecuaci\'{o}n (1b),
$S_A = f\,[1.25\,(S_c + S_h) - S_L]$, que es la aplicada por CAESAR II en el
presente modelo. En los reportes de CAESAR II, la columna \textit{Code}
corresponde al esfuerzo de c\'{o}digo calculado, la columna \textit{Allowable} al
admisible aplicable a cada caso, y el \textit{Ratio} a su cociente expresado en
porcentaje; un Ratio $\leq$ 100~\% implica cumplimiento. Para el material de la
l\'{i}nea (ASTM A312 TP 304L), el admisible b\'{a}sico de la Tabla A-1 de
B31.3-2016 es $S_c = S_h = 16.7$~ksi (\num{115.1}~\si{MPa}), valor consistente
con el admisible de expansi\'{o}n reportado por el modelo
(\num{269891.8}~\si{kPa}, secci\'{o}n~\ref{sec:resultados}).

% ─────────────────────────────────────────────────────────────────────────────
\subsection{Estructura del Documento}

La secci\'{o}n~\ref{sec:objetivos} presenta los objetivos del estudio; la
secci\'{o}n~\ref{sec:alcance} define su alcance y limitaciones; la
secci\'{o}n~\ref{sec:bases} re\'{u}ne las bases de dise\~{n}o congeladas del
proyecto; la secci\'{o}n~\ref{sec:metodologia} describe el flujo de trabajo y
los casos de carga; la secci\'{o}n~\ref{sec:resultados} presenta
desplazamientos, esfuerzos y cargas en soportes; la
secci\'{o}n~\ref{sec:analisis} desarrolla la evaluaci\'{o}n de cumplimiento de
c\'{o}digo; y las secciones~\ref{sec:conclusiones} y~\ref{sec:recomendaciones}
cierran con conclusiones y recomendaciones.
